\section{Problem Formulation and Hardness}
\label{sec:formulation}

\subsection{The redistricting graph model}

Let $G = (V, E, w_V, w_E)$ be a weighted planar graph where:
\begin{itemize}
  \item $|V| = n$ is the number of census tracts,
  \item $w_V : V \to \mathbb{Z}_{>0}$ assigns population to each tract,
  \item $w_E : E \to \mathbb{R}_{>0}$ assigns shared boundary length to
    each edge, and
  \item $G$ is planar (tracts tile a surface without crossings).
\end{itemize}

\begin{problem}[Balanced $k$-Cut Redistricting]
\label{prob:bkc}
Given $G = (V, E, w_V, w_E)$, integer $k \geq 2$, and tolerance
$\varepsilon > 0$, find a partition $\mathcal{P} = \{P_1, \ldots, P_k\}$
of $V$ that:
\begin{enumerate}
  \item minimises $\mathrm{EC}(\mathcal{P}) = \sum_{(u,v) \in E,\, \pi(u) \neq \pi(v)} w_E(u,v)$,
  \item satisfies $|w_V(P_i) - W/k| \leq \varepsilon W/k$ for all $i$
    (population balance), where $W = \sum_v w_V(v)$, and
  \item each $P_i$ induces a connected subgraph of $G$ (contiguity).
\end{enumerate}
\end{problem}

\subsection{NP-hardness of the redistricting problem}

\begin{theorem}[NP-hardness]
\label{thm:nphard}
Problem~\ref{prob:bkc} is NP-hard for planar graphs, even with $k = 2$
and unit vertex weights.
\end{theorem}

\begin{proof}[Proof sketch]
We reduce from the Connected Planar Graph Bisection problem, which asks
whether a planar graph can be bisected into two connected parts of equal
vertex count with edge cut at most $C$.
This problem is NP-complete by the results of Dyer and
Frieze~\citep{dyer1985partitioning} on connected graph partition on
planar graphs; see also Garey and Johnson~\citep{garey1979computers}
(Problem~ND14) for the disconnected variant, whose connected restriction
only strengthens hardness.

Given a Connected Planar Graph Bisection instance $(G', C)$ with $|V'|$
even, set all vertex weights $w_V \equiv 1$ (unit population), all edge
weights $w_E \equiv 1$ (unit boundary), $k = 2$, and $\varepsilon = 0$.
A valid solution to Problem~\ref{prob:bkc} on this instance consists of
two connected parts $P_1, P_2$ with $|P_1| = |P_2| = |V'|/2$ and total
edge cut minimised — which is precisely a minimum connected bisection of
$G'$.
Conversely, every minimum connected bisection of $G'$ yields a valid
solution to Problem~\ref{prob:bkc} (connectivity of parts satisfies
contiguity; equal sizes satisfy population balance with $\varepsilon = 0$).
The reduction is a polynomial-time many-one reduction since all
assignments are constant-time.
The planar graph structure of census tract adjacency graphs is preserved
under this reduction, as $G'$ is planar by assumption and weight
assignment does not alter planarity. \qed
\end{proof}

\noindent\textbf{Remark (feasibility vs.\ optimisation).}
Theorem~\ref{thm:nphard} establishes hardness of the \emph{optimisation}
variant (finding a partition with minimum edge cut).
The \emph{feasibility} variant --- does any valid partition exist? ---
is trivially solvable for connected graphs with $k \leq n$: any spanning
tree of $G$ can be partitioned into $k$ connected subtrees with
appropriate edge removal.
The NP-hardness result therefore concerns the difficulty of finding
an \emph{optimal} compact partition, not the existence of one.

\begin{proposition}[Approximation ratio]
\label{prop:approx}
No known polynomial-time algorithm achieves better than $O(\log n)$
approximation for balanced $k$-partition in general.
The approximation quality of METIS recursive bisection on redistricting-scale
planar graphs is empirically characterised in Section~\ref{sec:runtime};
no formal approximation guarantee for METIS's heuristic is claimed here.
\end{proposition}

\noindent\textbf{Context.}
The ARV framework~\citep{arora2009expander} achieves $O(\sqrt{\log n})$
approximation for Sparsest Cut via an SDP relaxation; this result applies
to Sparsest Cut, not to Balanced $k$-Cut with contiguity.
METIS's Fiduccia--Mattheyses (FM) local search~\citep{fiduccia1982linear}
has no formal connection to the ARV SDP bound.
By the Lipton--Tarjan separator theorem~\citep{lipton1979separator},
every planar graph has a balanced separator of size $O(\sqrt{n})$, which
gives an $O(\sqrt{n})$ edge-cut upper bound for a single bisection;
the approximation ratio this implies depends on the unknown optimum.
Section~\ref{sec:runtime} provides an empirical characterisation of how
closely METIS tracks the best partition found across multiple seeds.

\subsection{Empirical methodology}

We measure runtime by executing the METIS recursive bisection (C FFI,
version 5.1.0) on all 50 states across three census years (2000, 2010,
2020), with configurations:
\begin{itemize}
  \item \texttt{ufactor = 5}, \texttt{niter = 100}, \texttt{seed = 42}
  \item Single-threaded execution on a 16-core AMD Ryzen 9 5950X
  \item Runtime measured via \texttt{std::time::Instant} in the Rust
    harness (\texttt{redist-cli})
  \item Averaged over 10 independent runs to reduce variance
\end{itemize}

We fit $T = a \cdot n^b$ using ordinary least squares on $\log$-transformed
data, reporting $b \pm$ one standard error across the 50-state sample.
